\section{Future Work}
\label{ch:future_work}

The work presented in this thesis establishes a mathematically sound,
performant, and truly language-agnostic framework for architectural
dependency extraction. By abstracting the extraction logic through a
heuristic-based approach on Concrete Syntax Trees (CST) and
formalizing a flexible Intermediate Representation
($\mathcal{D}_{IR}$), the system successfully overcomes the
scalability and configuration bottlenecks of previous iterations. The
validation phase demonstrated the robustness of the model across
different object-oriented and procedural languages (such as Java,
Python, C++, and Rust), properly resolving complex nested scopes,
generic types, and cross-module imports.

However, as the primary focus of this research was the
conceptualization, formalization, and validation of the theoretical
model, several avenues remain open for future improvements and
extensions. The following sections outline the most promising
directions for future research.

\section{Formal Provability and Type Theory}

While the current framework provides an informal argument of
correctness by construction, achieved through explicit domains,
codomains, and well-typed functional composition ($\Phi = G \circ R
\circ E$), it does not currently employ a formal system of typing
judgments or inference rules (e.g., $\Gamma \vdash e : \tau$).
However, the architectural design of the system makes full
mathematical formalization a highly feasible and compelling direction
for future research.

The framework already possesses strong foundations suited for Type
Theory: the Intermediate Representation is built entirely on
Algebraic Data Types (sum and product types), and the resolution of
symbols follows a strict state machine (transitioning from
  \texttt{Unresolved} to \texttt{ResolutionQuery}, and finally to
\texttt{Resolved}).

To achieve complete formal provability, future iterations would need
to address several key gaps:
\begin{itemize}
  \item \textbf{Type-Level Phase Guarantees:} Currently, a
    \texttt{Module} is represented by the same Rust type both before
    and after resolution. The invariant that "no \texttt{Unresolved}
    references exist after Phase 2" is only guaranteed at runtime. A
    formal implementation would require parametric typing or
    Generalized Algebraic Data Types (GADTs), such as
    \texttt{Module<Phase>}, to enforce these state transitions
    strictly at compile time.
  \item \textbf{Formalization of Invariants:} To write rigorous
    inference rules (e.g., proving that if a module is well-formed in
      the unresolved state, the builder function will produce a
    well-formed query state), the structural invariants of what
    constitutes a "well-formed" module at each specific phase must be
    mathematically specified.
  \item \textbf{Conditional Correctness of Extraction:} Because Phase
    1 relies on heuristics, its extraction correctness is
    intrinsically approximate and cannot be formally proven for every
    possible malformed AST. Therefore, the formalization would focus
    on proving \textit{conditional correctness}: guaranteeing that
    \textit{if} Phase 1 produces a valid Intermediate Representation,
    \textit{then} the subsequent resolution phases are mathematically sound.
  \item \textbf{Mechanized Proofs:} Finally, fully validating this
    type-theoretical model would require transcribing the framework's
    definitions into a proof assistant, such as Coq, Lean, or Agda,
    to mechanically verify all preservation and progress lemmas.
\end{itemize}

\section{Extraction Phase: srcML Integration}

While the current heuristic-based extraction directly from
Tree-sitter CSTs eliminates the need for manual dispatch tables, it
remains the most delicate phase of the pipeline. The framework relies
on structural heuristics and string-matching of node names (e.g.,
identifying nodes containing \texttt{struct} or \texttt{class}).
Consequently, it is somewhat fragile to unexpected grammar updates or
drastic structural variations in Tree-sitter parsers.

A highly promising future development involves replacing or
augmenting the first extraction phase with an XML-based parser such
as \textbf{srcML}. As discussed in the Related Work chapter, srcML
standardizes source code into an XML format while preserving all
original text and structural context. Transitioning to srcML would
completely decouple the extraction engine from Tree-sitter's AST
volatility, transforming the heuristic identification of
architectural entities into a purely declarative and robust process
(e.g., using standardized XPath queries to reliably extract
\texttt{<class>} or \texttt{<function>} nodes).

\section{Performance Optimizations}

Throughout the development of this thesis, the primary objective was
ensuring the correctness and algorithmic scalability of the
mathematical model, the resolution strategies, and the algebraic
query engine. Consequently, there has been limited focus on low-level
execution optimization, and the current prototype operates almost
entirely sequentially. While the algorithmic complexity is
well-bounded, real-world industrial applications demand extremely
fast execution times to integrate seamlessly into Continuous
Integration (CI) pipelines.

Future iterations of the framework could achieve massive performance
gains by introducing concurrency at multiple levels of the pipeline.
The Extraction Phase (Phase 1) is "embarrassingly parallel": since
the Tree-sitter AST parsing and heuristic extraction do not rely on
any shared state, each source file can be processed entirely
independently. Implementing a thread-pool architecture (for instance,
utilizing Rust's \texttt{rayon} library) would allow the system to
process hundreds of files concurrently, bottlenecked only by the disk
I/O and available CPU cores.

The Resolution Phase (Phase 2), however, presents a more complex
concurrency challenge. The global Scope Tree and the Intermediate
Representation ($\mathcal{D}_{IR}$) act as shared memory structures.
Parallelizing the resolution process requires migrating from standard
sequential data structures to lock-free concurrent collections or
employing fine-grained Reader-Writer locks (e.g., \texttt{RwLock}).
For example, independent queries resolving variables inside different
function bodies can run perfectly in parallel by acquiring read-locks
on the global scope, while insertions of new entities or edges would
require isolated write-locks. Furthermore, batching graph insertions
instead of streaming them individually would drastically minimize
lock contention on the underlying graph database, enabling the tool
to analyze enterprise-scale, multi-million-line codebases in a
fraction of the current time.

\section{Inter-Language Dependency Detection}

Currently, while the framework successfully aggregates graphs
extracted from multi-language repositories into a unified
representation, the resolution of dependencies remains confined
within the boundaries of each individual language. A native Java
method call, for instance, is correctly identified as an external
invocation, but it is not formally linked to its actual C/C++ implementation.

In modern distributed or highly optimized systems, cross-language
interactions are ubiquitous. A compelling future extension involves
detecting and formally resolving these \textbf{inter-language
dependencies}. To achieve this natively, the resolution engine must
be extended with "bridge" heuristics capable of translating
cross-boundary calling conventions. For example, in the case of the
Java Native Interface (JNI), the framework could identify a
\texttt{native} Java method, automatically compute its C/C++ mangled
name (e.g., following the \texttt{Java\_Package\_Class\_Method}
convention), and dispatch a resolution query directly to the C/C++
namespace within the same IR. Similar bridging logic could be applied
to Foreign Function Interface (FFI) bindings between Python and Rust.

Beyond direct FFI, modern microservice architectures often
communicate via Remote Procedure Calls (RPC) defined by Interface
Definition Languages (IDLs), such as Protocol Buffers for gRPC or
OpenAPI specifications. By writing an extraction heuristic for these
interface files, the analyzer could use the IDL as a central
contract, automatically generating dependency edges between a Go
backend serving a request and a Python client invoking it. Extending
the resolution engine to bridge these language boundaries would allow
the extraction of a truly global architectural dependency graph. This
offers unprecedented visibility into complex polyglot architectures,
enabling system-wide refactoring and deep security auditing that
traditional single-language analyzers cannot provide.

\section{Alternative Programming Paradigms}

The empirical validation and heuristic design of the current
framework have primarily targeted procedural and object-oriented
programming paradigms (C-like languages, Java, Python, Rust). These
languages share foundational architectural concepts such as modules,
classes, structs, and explicitly invoked functions.

Although the Intermediate Representation ($\mathcal{D}_{IR}$) was
designed mathematically to be as agnostic as possible, extending the
framework to support entirely different programming paradigms
represents a significant future challenge. Future research should
evaluate and expand the model's capability to capture architectural
dependencies in \textbf{Functional} languages (e.g., Lisp, Haskell,
Clojure), which heavily rely on higher-order functions and closures,
or in \textbf{Logic} programming languages (e.g., Prolog).
Incorporating these paradigms will likely require introducing new
heuristic classifiers in the extraction phase and refining the query
algebra to accommodate exotic scoping rules and
pattern-matching-based execution semantics.
